\documentclass[11pt,a4paper]{article}
\usepackage[T1]{fontenc}
\usepackage[utf8]{inputenc}
\usepackage{lmodern}
\usepackage{amsmath,amssymb,amsthm,mathtools}
\usepackage{geometry}
\geometry{margin=2.5cm}

\title{Backpropagation from First Principles}
\author{nabla-flow}
\date{}

\begin{document}
\maketitle

\section{Scope}
This document derives reverse-mode automatic differentiation (backpropagation)
from first principles for scalar loss functions
\(L: \mathbb{R}^n \to \mathbb{R}\).

\section{Derivative and partial derivative from limits}
For \(f: \mathbb{R} \to \mathbb{R}\), the derivative at \(x\) is
\begin{equation}
  f'(x) = \lim_{h \to 0} \frac{f(x+h)-f(x)}{h}.
\end{equation}

For \(f: \mathbb{R}^n \to \mathbb{R}\), the \(i\)-th partial derivative is
\begin{equation}
  \frac{\partial f}{\partial x_i}(x)
  = \lim_{h \to 0}
  \frac{f(x_1,\dots,x_i+h,\dots,x_n)-f(x_1,\dots,x_n)}{h}.
\end{equation}

The gradient is
\begin{equation}
  \nabla f(x)
  = \left(\frac{\partial f}{\partial x_1}(x),\dots,\frac{\partial f}{\partial x_n}(x)\right)^\top.
\end{equation}

\section{Computational graph and reverse accumulation}
Assume a directed acyclic graph (DAG) with intermediate variables
\(v_1,\dots,v_m\) and scalar output \(L\).
Define the adjoint
\begin{equation}
  \bar v := \frac{\partial L}{\partial v}.
\end{equation}

If node \(u\) has direct successors \(w_1,\dots,w_r\), then by chain rule
\begin{equation}
  \bar u
  = \frac{\partial L}{\partial u}
  = \sum_{k=1}^{r} \frac{\partial L}{\partial w_k}\frac{\partial w_k}{\partial u}
  = \sum_{k=1}^{r} \bar w_k\,\frac{\partial w_k}{\partial u}.
\end{equation}

This identity is the reverse accumulation rule.

\section{Local derivatives from first principles}
\subsection{Addition}
Let \(z = x+y\).
\begin{equation}
  \frac{\partial z}{\partial x}
  = \lim_{h\to0}\frac{(x+h)+y-(x+y)}{h} = 1,
  \qquad
  \frac{\partial z}{\partial y}=1.
\end{equation}
Backprop update:
\begin{equation}
  \bar x \mathrel{+}= \bar z,
  \qquad
  \bar y \mathrel{+}= \bar z.
\end{equation}

\subsection{Multiplication}
Let \(z = xy\).
\begin{equation}
  \frac{\partial z}{\partial x}
  = \lim_{h\to0}\frac{(x+h)y-xy}{h} = y,
  \qquad
  \frac{\partial z}{\partial y}=x.
\end{equation}
Backprop update:
\begin{equation}
  \bar x \mathrel{+}= \bar z\,y,
  \qquad
  \bar y \mathrel{+}= \bar z\,x.
\end{equation}

\subsection{Power with integer exponent}
Let \(z=x^n\) with \(n\in\mathbb{N}\). Using the binomial theorem:
\begin{align}
  \frac{(x+h)^n-x^n}{h}
  &= \frac{\sum_{k=0}^{n}\binom{n}{k}x^{n-k}h^k - x^n}{h} \\
  &= nx^{n-1} + \sum_{k=2}^{n}\binom{n}{k}x^{n-k}h^{k-1}.
\end{align}
Taking \(h\to0\) eliminates all terms \(k\ge2\), so
\begin{equation}
  \frac{\partial z}{\partial x}=nx^{n-1},
  \qquad
  \bar x \mathrel{+}= \bar z\,nx^{n-1}.
\end{equation}

\section{Worked chain-rule example}
Let
\begin{equation}
  L=(ab+c)d.
\end{equation}
Introduce intermediates
\begin{equation}
  u=ab,
  \qquad
  v=u+c,
  \qquad
  L=vd.
\end{equation}

Initialize \(\bar L=1\).
\begin{align}
  \bar v &= \bar L\frac{\partial L}{\partial v}=d, &
  \bar d &= \bar L\frac{\partial L}{\partial d}=v, \\
  \bar u &= \bar v\frac{\partial v}{\partial u}=d, &
  \bar c &= \bar v\frac{\partial v}{\partial c}=d, \\
  \bar a &= \bar u\frac{\partial u}{\partial a}=bd, &
  \bar b &= \bar u\frac{\partial u}{\partial b}=ad.
\end{align}
Therefore
\begin{equation}
  \frac{\partial L}{\partial a}=bd,
  \quad
  \frac{\partial L}{\partial b}=ad,
  \quad
  \frac{\partial L}{\partial c}=d,
  \quad
  \frac{\partial L}{\partial d}=ab+c.
\end{equation}

\section{Finite-difference gradient check}
For implementation validation only (not training), use central differences:
\begin{equation}
  \frac{\partial L}{\partial x_i}
  \approx
  \frac{L(x_i+\varepsilon)-L(x_i-\varepsilon)}{2\varepsilon},
  \qquad
  0<\varepsilon\ll1.
\end{equation}

\section{References}
\begin{itemize}
  \item M. Spivak, \textit{Calculus}.
  \item W. Rudin, \textit{Principles of Mathematical Analysis}.
  \item C. M. Bishop, \textit{Pattern Recognition and Machine Learning}.
  \item I. Goodfellow, Y. Bengio, A. Courville, \textit{Deep Learning}.
  \item J. Nocedal, S. J. Wright, \textit{Numerical Optimization}.
\end{itemize}

\end{document}
